\bookchapter{第四十章：需求工程——把想法變成AI可執行的指令}

\chapterintro{如果說上下文工程解決的是“讓 AI 瞭解背景”，那需求工程解決的是“讓 AI 知道這次具體要做什麼”。這是 AI 程式設計中最能體現人與人差距的環節——同樣的工具、同樣的模型，有人能做出令人驚豔的專案，有人只能做出平庸的玩具。區別往往不在技術能力，而在需求表達能力。}

\booksection{一、從模糊想法到結構化需求}

每個專案都始於一個模糊的想法——“我想做一個記帳 App”“我想做個個人網站”“我想做個待辦清單工具”。這些想法本身沒有問題，但如果你直接把它丟給 AI，產出的品質完全取決於運氣。需求工程的目標，就是把這個模糊的想法，變成一個 AI 能精確執行的結構化需求。

好訊息是：結構化需求不需要你寫一份幾十頁的 PRD。它只需要你回答五個問題，我們稱之為“需求五要素”：

\begin{center}
\begin{tabularx}{\linewidth}{|X|X|X|}
\hline
\textbf{要素} & \textbf{問題} & \textbf{示例：記帳 App} \tabularnewline
\hline
Who & 給誰用？ & 個人使用，記錄日常收支 \tabularnewline
What & 做什麼？ & 一個能記錄每筆收支、按月彙總、視覺化的應用 \tabularnewline
Why & 為什麼做？ & 想了解自己的消費結構，找出不必要的支出 \tabularnewline
How & 大致怎麼做？ & 前端用 React，資料存瀏覽器本地，先不做後端和登入 \tabularnewline
Constraint & 有什麼限制？ & 單人開發，2 周內完成，不需要上線部署 \tabularnewline
\hline
\end{tabularx}
\end{center}

\begin{center}\small\color{gray}需求五要素及示例\end{center}

這五個要素不需要寫成正式檔案，甚至在腦海中快速過一遍就足夠了。關鍵是養成“先想清楚再開口”的習慣。你會發現，當你能清晰地回答這五個問題時，你的需求描述自然就變得結構化了。

\booksection{二、需求拆解的分層藝術}

對於較大的專案，一個需求描述是不夠的。你需要把大需求拆解成小的、可獨立驗證的子需求。這個過程叫做“需求拆解”，核心原則是：每個子需求都應該能獨立完成、獨立驗證、獨立交付。

以“做一個個人記帳 App”為例，一個合理的拆解方式是：

\begin{enumerate}
  \item 實現基礎的收支記錄功能（能新增、檢視、刪除記錄）；
  \item 新增分類功能（給每筆記錄打上餐飲、交通等標籤）；
  \item 實現月度彙總功能（按月統計收支情況）；
  \item 新增資料視覺化（用圖表展示消費結構）。
\end{enumerate}

每一步都是一個可以獨立執行和驗證的功能增量。

需求拆解的黃金法則：每一步完成後，專案都應該是一個"能用的東西"——哪怕功能很簡陋。永遠不要讓專案停留在"什麼都做了一半但什麼都用不了"的狀態。

這種拆解方式有兩個好處。首先，它讓你始終能看到進展——每完成一步，你的 App 就多了一個功能，這種正反饋對於長期專案至關重要（人都是需要成就感的嘛）。其次，它讓 AI 的每次任務都足夠小，降低了“AI 理解偏差”的風險——任務越小，AI 的理解越準確，產出品質越高。

\booksection{三、PRD 寫作模板與實踐}

對於較大的專案（比如一個有前後端、多個功能模組的全棧應用），僅憑需求五要素可能不夠。這時你需要寫一份更正式的專案需求檔案（PRD）。一個容易誤解的點是：PRD 不是寫給 AI 看的——它是寫給你自己的，幫你理清思路，然後你才能給 AI 提供更精確的指令。

一份好的 PRD 不需要很長，但應該包含六個部分：專案概述（一段話說清楚做什麼）、功能列表（每個功能點是什麼）、技術方案（用什麼技術棧、大致架構）、資料模型（核心資料結構）、介面描述（主要頁面長什麼樣）、開發計劃（分幾步完成）。

\begin{lstlisting}
# 個人記帳應用 PRD
## 專案概述
一個輕量級個人記帳應用，支援收支記錄、分類、月度彙總和資料視覺化。
## 功能列表
1. 記錄收支：金額、型別(收入/支出)、分類、備註、日期
2. 分類管理：預設餐飲、交通、購物、娛樂、其他
3. 月度彙總：按月統計總收入、總支出、結餘
4. 消費視覺化：餅圖展示各分類佔比
5. 資料匯出：匯出為CSV
## 技術方案
- 前端：React + TypeScript + Tailwind CSS
- 資料儲存：localStorage（無需後端）
- 圖表：使用Chart.js
- 構建：Vite
## 資料模型
Record {
  id: string
  amount: number
  type: 'income' | 'expense'
  category: string
  note: string
  date: string  // YYYY-MM-DD
}
## 開發計劃
Step 1: 基礎記錄功能 (CRUD)
Step 2: 分類管理
Step 3: 月度彙總
Step 4: 資料視覺化
Step 5: 資料匯出
\end{lstlisting}

注意這份 PRD 的寫作風格：用簡潔的列表和結構化的格式，而不是長篇大論。每個功能點都有明確的定義，資料模型用簡單的型別描述。這種風格的 PRD 不僅你自己看得清楚，當你把它提供給 AI 時，AI 也能精確理解每個需求點。

\booksection{四、實戰：從想法到方案的完整推演}

紙上談兵沒意思，咱們用一個完整的例子走一遍需求工程的全過程。假設你的想法是：“我想做一個能追蹤我每天學了什麼的學習追蹤器”。

第一步，過一遍需求五要素。

\begin{itemize}
  \item Who：個人使用。
  \item What：一個能記錄每天學習內容、追蹤學習時長、回顧學習進度的工具。
  \item Why：想養成持續學習的習慣，看到自己的積累。
  \item How：先做一個網頁版，資料存在瀏覽器本地。
  \item Constraint：零基礎開發，1 周內完成核心功能。
\end{itemize}

第二步，拆解需求。

\begin{enumerate}
  \item Step 1：實現"記錄學習條目"功能（日期、內容、時長）；
  \item Step 2：新增"學習時長統計"（今日、本週累計）；
  \item Step 3：實現"學習日曆"檢視（像 GitHub 貢獻圖一樣展示）；
  \item Step 4：新增"學習回顧"功能（隨機展示過去某天的學習記錄，幫助複習）。
\end{enumerate}

第三步，寫一份簡短 PRD。不需要長——上面的拆解本身就構成了 PRD 的主體。把它整理成一個 .md 檔案，放在專案根目錄，作為你和 AI 的共同參考。

第四步，開始第一步的開發。給 AI 的指令可以是：

\begin{lstlisting}
我需要一個學習追蹤器的第一個功能——記錄學習條目。
使用者可以新增一條學習記錄，包含日期、學習內容、學習時長（分鐘）。
資料存在localStorage。
前端用React + TypeScript。
請給出完整實現。
\end{lstlisting}

注意這個指令的特點：它只描述當前這一步的目標（記錄功能），不摻雜其他功能；它明確了技術約束（React、TypeScript、localStorage）；它給了足夠的資訊讓 AI 知道“做什麼”和“怎麼做”，但具體實現由 AI 決定。這就是一個經過需求工程加工後的、高品質的 AI 程式設計指令。

\begin{tcbpractice}{來動手：給你的"心頭好"寫一份五要素卡}
選一個你醞釀已久的小專案，花 5 分鐘填一張卡：Who / What / Why / How / Constraint，各一句話。然後把它拆成 3-5 個“每一步都能用”的步驟。你會發現：把想法寫下來之後，它就從“白日夢”變成了“待辦清單”。

\end{tcbpractice}

\begin{tcbsummary}{本章小結}
\begin{itemize}
  \item 需求五要素：Who / What / Why / How / Constraint——想清楚再開口
  \item 需求拆解黃金法則：每一步完成後專案都是"能用的東西"
  \item PRD 是寫給自己的：概述、功能列表、技術方案、資料模型、介面、計劃
  \item 好指令 = 只描述當前一步 + 明確約束 + 說"做什麼"不說"怎麼做"
  \item 完整推演模板：五要素 拆解 PRD 給指令，四步走
\end{itemize}

\end{tcbsummary}
